The central result is not that one prostate-cancer target explains a foundation model. Six
clinically interpretable axes occupied distinct positions in the evidence map: Grade/ISUP and
tumor phenotype/content had the broadest representation-level transport; PTEN-related
information was recoverable and AR activity aligned positively, but both became conditional;
SPOP did not yield a qualified direction; and recurrence changed with endpoint, comparator,
representation, and scale. The ISUP intervention extended one eligible axis from representation
availability to internal functional sensitivity. Site-heldout evidence was encoder-specific and
the locked LEOPARD test did not qualify external transport, so ISUP does not become a sole or
externally stable organizing target.

Taken together, these contrasts---rather than any single performance value---support the paper's
claim of selective, evidence-qualified interpretability. The hierarchy turns familiar targets
into contestable coordinates rather than post-hoc labels attached to an AI output. Transported coordinates can anchor clinician
review of agreement and disagreement; conditional or unsupported coordinates indicate where
an explanation should be withheld; and missing evidence identifies the next cohort,
intervention, or endpoint to test. This supports reliability auditing in clinically intelligible
terms, but whether presenting the map improves calibrated reliance or decisions remains a
prospective question.

Reference availability constrains the hierarchy. The study did not ask a slide-only model to
agree with six unobserved quantities: each estimate used an available paired reference. Grade
and phenotype had the widest compatible external coverage, molecular references were largely
confined to TCGA-PRAD, and recurrence definitions were non-equivalent. Missing external or
same-patient functional truth therefore produced not-evaluated or blocked states, whereas SPOP
was unsupported after comparison with an available genomic reference. The hierarchy reflects
both biological accessibility in H\&E and unequal opportunities for transport and functional
validation.

The target-specific results reinforce this distinction. Morphology-linked signals transported
most clearly, although PANDA phenotype used grade-derived tumor status rather than exact tumor
fraction. PTEN and AR support conditional histology--molecular associations, not molecular
diagnosis from H\&E or BCR-head use. Recurrence carried endpoint-associated risk information,
but its variation across event definitions and clinical references precludes a validated
prognostic-biomarker claim. The correlated setting grid exposed analytic sensitivity without
creating new patients, institutions, or independent replications.

Each result therefore has a different argumentative role. Morphology provides the positive
cross-cohort audit anchors; PTEN and AR demonstrate why decodability must be qualified by
conditional and site evidence; SPOP shows that an available reference can yield a principled
rejection rather than an invented explanation; recurrence shows that the endpoint and comparator
are part of the alignment claim; and ISUP alone supplies an internal link to a locked judgment.

Our approach also differs from methods that directly connect concepts to a trained prediction.
TCAV estimates prediction sensitivity along a concept direction, and concept bottleneck models
route predictions through explicit concepts \cite{kim2018,koh2020}. Related medical-imaging and
pathology systems connect influential regions or embeddings to recognizable findings
\cite{lou2026,han2026}. Here, the primary contribution is comparative qualification of
pre-existing targets in frozen representations. Targeted erasure is a secondary worked example,
not a complete concept-grounded architecture or mechanistic-fidelity result.

Internal ISUP erasure exceeded matched-random removal in both encoders, but refit intervals did
not establish indispensability and the limited ISUP-only comparison was not a comprehensive
clinical-increment study. Only
Virchow passed site-heldout TCGA. In LEOPARD, CONCH erasure was positive but its head interval
crossed chance; Virchow showed neither valid discrimination nor targeted sensitivity. Thus the
result is internal subspace sensitivity plus a failed-or-inconclusive external test, not a
human-equivalent mechanism. Exact rerun agreement does not rescue the failed external gate. Absence of a
functional-use result is therefore not one common negative result: other axes stopped for
missing truth, unrun protocols, an unqualified direction, or a non-applicable question.

Several limitations remain. Morphology had greater external-reference availability; some rows
lacked intervals or event accounting; site was confounded with acquisition and case mix; and
recurrence analyses were post hoc or event-limited. Failed independent-domain detector sensitivity keeps the
intervention whole-tissue. LEOPARD lacked ISUP and treatment covariates, may use a non-equivalent
BCR definition, and had prior outcome access; it was not prospectively untouched. Calibration,
treatment utility, autonomous diagnosis, population-wide generalization, clinician trust,
reliance, diagnostic accuracy, and patient outcomes were not evaluated. Familiar explanations
can themselves increase inappropriate reliance \cite{prinster2024}, making prospective testing
essential.

Finally, the map creates a disciplined discovery space rather than reporting a new biomarker.
After reliable known coordinates and technical confounders are jointly accounted for, a residual
signal should advance only if it recurs across models and independent cohorts, contributes to a
locked prediction head, can be defined as a reproducible human-measurable feature, and receives
biological or clinical validation. This separates potential biomarker discovery from artifacts
while preserving the present contribution: a qualified map of what is shared, conditional,
unsupported, and still untested.
